\documentclass[11pt]{article}
\usepackage{fontspec}
\usepackage[french]{babel}
\usepackage{amsmath}
\usepackage{amssymb}
\usepackage{amsthm}
\usepackage{mathrsfs}
\usepackage{tikz-cd}
\usepackage{enumitem}

\newcommand{\printedpage}[1]{\marginpar{\footnotesize p.~#1}}
\newcommand{\msp}[1]{\marginpar{\footnotesize [#1]}}
\newcommand{\Ob}{\operatorname{Ob}}
\newcommand{\Hom}{\operatorname{Hom}}
\newcommand{\Ind}{\operatorname{Ind}}
\newcommand{\card}{\operatorname{card}}
\newcommand{\cM}{\mathcal{M}}
\newcommand{\cE}{\mathcal{E}}
\newcommand{\cC}{\mathcal{C}}
\newcommand{\cQ}{\mathcal{Q}}
\newcommand{\uS}{\underline{S}}
\newcommand{\uT}{\underline{T}}
\newcommand{\U}{\mathfrak{U}}

\begin{document}
\setcounter{footnote}{195}

\printedpage{184}
% head: A. Grothendieck : Dérivateurs, Ch. XVIII, Édition des Universités de Montpellier II et Paris VII
% foot: Version 17 décembre 2025

\msp{page 475}
On trouve alors le

\begingroup\itshape
\noindent\textbf{\upshape Corollaire 4.16.6.2.} Soit $f : \cE \longrightarrow \cC$ un foncteur. Conditions équivalentes :
\begin{enumerate}[label={}]
\item[a)] Pour toute catégorie $\cC'$ \textup{petite} et \textup{karoubienne}, et tout foncteur $\cC' \longrightarrow \cC$, $\cE' = \cE \times_{\cC} \cC'$ est accessible, et le foncteur $\cE' \longrightarrow \cC'$ est accessible.
\item[b)] Comme a), mais $\cC'$ étant une des trois catégories $\Delta_1$, $\cQ$, $\cQ'$ (où $\cQ$, $\cQ'$ sont définies dans 4.15.6.1 ci-dessus).
\item[b$'$)] \textup{(}\footnote{?, à vérifier avec soin !}\textup{)}. (Si $\cC$ est sans projecteurs.) Comme b), mais en se limitant au cas $\cC' = \boldsymbol{\Delta}_1$ (à l'exclusion de $\cQ$, $\cQ'$).
\end{enumerate}
\endgroup

\begingroup\itshape
\noindent\textbf{\upshape Corollaire 4.16.6.3.} Supposons $f : \cE \longrightarrow \cC$ fibrant ou cofibrant. Alors les conditions équivalentes précédentes impliquent que :
\begin{enumerate}[label={}]
\item[1$^\circ$)] Les $\cE_s$ sont accessibles ($s \in \Ob \cC$).
\item[2$^\circ$)] Pour toute flèche $\lambda : s \longrightarrow t$ dans $\cC$, le foncteur $\lambda^* : \cE_t \longrightarrow \cE_s$ (resp. $\cE_s \longrightarrow \cE_t$) est accessible.
\end{enumerate}
\noindent D'autre part, l'ensemble de ces conditions 1$^\circ$, 2$^\circ$) équivaut à la validité de b) ci-dessus, dans le cas $\cC' = \boldsymbol{\Delta}_1$, ou encore à la validité
\msp{page 476}
de a) ci-dessus, quand $\cC'$ est \textup{sans projecteurs}.
\endgroup

\textbf{Question 4.16.6.4.} \textit{Les conditions 1$^\circ$ et 2$^\circ$ précédentes impliquent-elles celles de 4.16.6.2 ? En d'autres termes, si la catégorie base $\cC$ est égale à $\cQ$ ou à $\cQ'$, alors ces conditions impliquent-elles que $\cE$ est accessible, et $f : \cE \longrightarrow \cC$ accessible ?}

Ça semble idiot de ne pas connaître la réponse !

\textbf{4.16.6.5.} J'ai envie d'appeler catégorie \textit{accessible sur $\cC$}, une catégorie $\cE$ sur $\cC$ satisfaisant les conditions équivalentes de 4.16.6.1 -- du moins dans le cas où on suppose $\cC$ petite. Pour que cette terminologie soit vraiment satisfaisant, il faudrait que l'on sache prouver que ces conditions impliquent ceci : pour toute catégorie accessible $\cM$ et tout foncteur accessible $\cM \longrightarrow \cC$, la catégorie $\cE \times_{\cC} \cM$ est accessible, et le foncteur $\cE \times_{\cC} \cM \longrightarrow \cM$ est accessible. % [sic] satisfaisant -> satisfaisante, accord avec « terminologie »

\msp{page 476 bis}
L'idée, c'est qu'on appellerait (sous hypothèse de petitesse) \textit{catégorie accessible $\cE$} sur une autre $\cC$, une catégorie ayant la propriété qu'on vient de dire -- mais en supposant maintenant de plus que $\cC$ satisfasse à $\mathrm{L}_\pi$, j'imagine, pour tout foncteur accessible $\cM \longrightarrow \cC$, avec $\cM$ accessible, $\cE \times_{\cC} \cM$ est accessible et le foncteur canonique $\cE \times_{\cC} \cM \longrightarrow \cE$ est accessible.

\msp{page 477}
\subsection*{4.16.7. Retour sur les sous-catégories fermées d'une catégorie accessible (4.16.1).}

\textbf{Proposition 4.16.7.1.} \textit{Soient $\cM$ une catégorie accessible, $(\uS_i)_{i\in I}$ une famille de sous-catégories fermées accessibles dans $\cM$. On ne suppose que \textup{card}\,$I$ soit petit. Soit $\uS$ la sous-catégorie strictement pleine de $\cM$ telle que $\Ob \uS = \bigcup_{i\in I} \Ob \uS_i$. Alors pour que $\uS$ % [sic] On ne suppose que -> On ne suppose pas que, « pas » apparemment omis
\printedpage{185}
% head: A. Grothendieck : Dérivateurs, Ch. XVIII, Édition des Universités de Montpellier II et Paris VII
% foot: Version 17 décembre 2025
soit accessible dans $\cM$, il faut et il suffit qu'il existe une \textup{petite} partie $J \subset I$, telle que $\Ob \uS = \bigcup_{i\in J} \uS_i$. (En particulier, si $I$ est ordonnée, et $i \longmapsto \uS_i$ croissante et grande devant tout petit cardinal, il faut et il suffit qu'il existe un $i \in I$ tel que $\uS = \uS_i$, i.e. tel que $j \geqslant i$ implique $\uS_j = \uS_i$.)}

C'est évidemment suffisant, en vertu de 4.13\dots{} Prouvons inversement que si $\uS$ est accessible dans $\cM$, alors la condition précédente est satisfaite. Soit $\pi$ un cardinal $\geqslant \pi_0$ tel que $\uS$ soit $\pi$-accessible dans $\cM$, donc
\[
\uS \ \xleftarrow{\ \sim\ }\ \Ind_\pi \uS(\pi) \ , \qquad\qquad \text{où } \uS(\pi) = \uS \cap \cM_\pi.
\]

\msp{page 478}
Soit $\uS(\pi)_0$ une sous-catégorie réduite de $\uS(\pi)$, telle que $\uS(\pi)_0 \hookrightarrow \uS(\pi)$ soit une équivalence, de sorte que $\uS(\pi)_0$ est petite. Pour tout $x \in \Ob \uS(\pi)_0$, soit $i_x \in I$ tel que $x \in \Ob \uS_{i_x}$. Soit $J$ l'ensemble de tous les $i_x$, que est une petite partie de $I$, et soit $\uS_J$ la sous-catégorie strictement pleine de $\cM$ ayant pour ensemble d'objets $\bigcup_{j\in J} \Ob \uS_j$. Donc $\uS(\pi)_0$ est contenue dans $\uS_J$, donc aussi $\uS(\pi)$, puisque $\uS_J$ est strictement pleine : % [sic] que est -> qui est, pronom relatif sujet
\[
\uS(\pi) \ \hookrightarrow\ \uS_J \ .
\]
D'ailleurs $\uS_J$ est accessible dans $\cM$ (loc. cit.), et est fermée dans $\cM$ comme réunion de sous-catégories fermées, \textit{a fortiori elle satisfait à $L_\pi$}. Donc, comme $\uS(\pi) \subset \uS_J$, le foncteur essentiellement surjectif $\Ind_\pi(\uS(\pi)) \longrightarrow \uS$ se factorise par $\uS_J$. Donc $\uS_J \hookrightarrow \uS$ est essentiellement surjectif, et comme $\uS_J$ est strictement pleine, on en conclut $\uS_J = \uS$, q.e.d.

\msp{page 479}
\textbf{Proposition 4.16.7.2.} \textit{Soient $\cM$ une catégorie accessible, $x$ un objet de $\cM$, et ${}_x\cM$ la sous-catégorie fermée de $\cM$ engendrée par $x$, formée des $y \in \Ob \cM$ tels que $\Hom(x,y) \neq \emptyset$, i.e. l'image essentielle de $x\backslash\cM \longrightarrow \cM$. Alors ${}_x\cM$ est une sous-catégorie de $\cM$ accessible dans $\cM$.}

En effet, c'est l'image inverse de la sous-catégorie $\mathrm{Ens}^*$ accessible dans $\mathrm{Ens}$ par le foncteur représentable (donc accessible)
\[
h_x \ :\ y \longmapsto \Hom(x,y) \ :\ \cM \ \longrightarrow\ \mathrm{Ens} \ ,
\]
donc est accessible dans $\cM$ en vertu de 4.13\dots

\textbf{Corollaire 4.16.7.3.} \textit{Soit $\cM$ une catégorie accessible, $\uS$ une sous-catégorie fermée. Pour que $\uS$ soit accessible, il faut et il suffit qu'il existe une \textup{petite} partie $J$ de $\uS$, telle que $\Ob \uS = \bigcup_{x\in J} \Ob\,(\,{}_x\cM)$.}

En effet, on a
\[
\Ob \uS \ =\ \bigcup_{x\in\Ob\uS} \Ob\,(\,{}_x\cM) \ ,
\]
i.e. $\uS$ apparaît comme réunion de
\msp{page 480}
la famille (pas nécessairement petite) de sous-catégories fermées ${}_x\cM$ de $\cM$ (indexées par $x \in \Ob \uS$), accessibles dans $\cM$ en vertu de 4.16.7.2. Donc notre assertion est un cas particulier de 4.16.7.1.

\printedpage{186}
% head: A. Grothendieck : Dérivateurs, Ch. XVIII, Édition des Universités de Montpellier II et Paris VII
% foot: Version 17 décembre 2025

On peut aussi, pour toute partie petite $J$ de $\Ob \uS$, considérer la sous-catégorie fermée $\uS_J \subset \uS$ de $\cM$ engendrée par $J$. L'ensemble $I$ des parties $J$ en question, ordonné par inclusion, est grand devant tout cardinal $\pi \in \U$, et $\uS$ est réunion croissante des $\uS_J$. D'autre part, la proposition nous dit que $\uS$ est accessible dans $\cM$ si et seulement s'il existe un $J \in I$, tel que $\uS_J = \uS$, i.e. tel que $\uS_J = \uS_{J'}$ pour $J \subset J'$.

\begingroup\itshape
\noindent\textbf{\upshape Proposition 4.16.7.4.} Soit $\cM$ une catégorie accessible. Les conditions suivantes sont équivalentes :
\begin{enumerate}[label={}]
\item[a)] Il existe une sous-catégorie $\uS$ de $\cM$ fermée dans $\cM$, et non accessible dans $\cM$.
\item[b)] Soit $I$ l'ensemble des ordinaux $\alpha$ tels que $\card \alpha \in \U$, i.e. des
\msp{page 481}
ordinaux $\alpha$ qui sont dans $\U$. Donc $I = I_{\alpha_0}$, où $\alpha_0$ est le plus petit ordinal qui n'appartienne pas à $\U$, et $I$ est grand devant tout petit cardinal $\pi$, i.e. devant tout cardinal dans $\U$ \textup{(}\footnote{introduire $I = I_{\U}$ avant l'énoncé de la proposition.}\textup{)}. La condition envisagée est la suivante :

Il existe une suite transfini % [sic] transfini -> transfinie, accord avec « suite »
\[
\underline{x} \ =\ (x_\alpha)_{\alpha\in I}
\]
d'éléments de $\Ob \cM$, telle que l'on ait
\[
\Hom(x_\beta, x_\alpha) \ =\ \emptyset \qquad\qquad \textit{si } \beta < \alpha.
\]
De façon précise, pour une telle suite transfinie, soit $\uS$ ou ${}_{\underline{x}}\cM$ la sous-catégorie fermée de $\cM$ engendrée par la famille des $x_\alpha$, i.e. telle que $\Ob \uS = \bigcup_{\alpha\in I} \Ob\ {}_{x_\alpha}\cM = \bigcup_{\alpha\in I} x_\alpha\backslash\Ob \cM$. Alors ${}_{\underline{x}}\cM$ est une sous-catégorie fermée non accessible dans $\cM$, et toute sous-catégorie fermée non accessible dans $\cM$ s'obtient de cette façon.
\end{enumerate}
\endgroup

\textsc{Démonstration}. Considérons une suite $\underline{x} = (x_\alpha)_{\alpha\in I}$ comme dessus. Si $J$ est une petite partie de $I$, il existe un $\alpha \in I$ qui majore les $\beta \in J$, donc $\Hom(x_\beta, x_\alpha) = \emptyset$, pour $\beta \in J$,
\msp{page 482}
ce qui signifie que $x_\alpha$ n'appartient à aucune des $\uS_\beta = {}_{x_\beta}\cM$, pour $\beta \in J$. Donc $\uS$ n'est par réunion d'une petite sous-famille de la famille des ${}_{x_\alpha}\cM$, donc en vertu de 4.16.7.1, $\uS$ n'est par accessible dans $\cM$. % [sic] n'est par réunion -> n'est pas réunion ; n'est par accessible -> n'est pas accessible, « par » pour « pas » (2 fois)

Inversement, soit $\uS$ une telle sous-catégorie dans $\cM$. Considérons un bon ordre sur $\Ob \uS$ (objets $s_\alpha$), et construisons la suite des $x_\alpha$ ($\geqslant s_\alpha$) par récurrence transfinie (\footnote{\textbf{N.B.} On a $\card \uS \notin \U$, sinon, $\uS$ serait petite, donc accessible dans $\cM$, contrairement à l'hypothèse.}).

1) On choisit $x_0 =$ premier élément de $\Ob \uS$ (existe, car $\Ob \uS \neq \emptyset$, i.e. $\uS \neq \emptyset$, car sinon, $\uS$ serait accessible dans $\cM$).

2) Si les $x_\beta$ ($\beta < \alpha$) sont construits, on prend pour $x_\alpha$ le premier élément de $\uS$ majorant les $x_\beta$ et qui n'appartient pas à la sous-catégorie fermée de $\cM$ engendrée par les $x_\beta$ ($\beta < \alpha$). (Cette sous-catégorie fermée est $\subset \uS$, et $\neq \uS$, car $\uS$ n'est pas accessible dans $\cM$.) On a donc bien
\[
\Hom(x_\beta, x_\alpha) \ =\ \emptyset \qquad\qquad \text{si } \beta < \alpha.
\]

\printedpage{187}
% head: A. Grothendieck : Dérivateurs, Ch. XVIII, Édition des Universités de Montpellier II et Paris VII
% foot: Version 17 décembre 2025

\msp{page 483}
Ainsi, la récurrence transfinie se poursuit sur $I$ tout entier. On voit par récurrence transfinie qu'on a, pour tout $\alpha \in I$,
\[
s_\alpha \ \leqslant\ x_\alpha
\]
(card $x_\alpha > x_\beta \geqslant s_\beta$ pour $\beta < \alpha$, donc $x_\alpha > s_\beta$ pour tout $\beta < \alpha$, ce qui équivaut à $x_\alpha \geqslant s_\alpha$). De plus, on trouve aussi, par récurrence transfinie, que
\[
\beta \ <\ \alpha \qquad\quad \Longrightarrow \qquad\quad
\text{\parbox[c]{5.5cm}{$s_\beta$ appartient à la sous-catégorie fermée $\uS_\alpha$ de $\cM$ engendrée par les $x_{\beta'}$ avec $\beta' < \alpha$.}}
\quad ,
\tag*{$(*)$}
\]
puisque par construction $s_\alpha$ est le \textit{premier} des éléments de $\uS$ qui n'est \textit{pas} dans ladite sous-catégorie $\uS_\alpha$ de $\cM$. Cela montre que tout $s_\beta$ appartient à la sous-catégorie fermée ${}_{\underline{x}}\cM$ de $\cM$ engendrée par la famille de \textit{tous} les $x_\alpha$, donc on trouve que $\uS \subset {}_{\underline{x}}\cM \subset \uS$, donc $\uS = {}_{\underline{x}}\cM$, q.e.d.

\begingroup\itshape
\noindent\textbf{\upshape Proposition 4.16.7.5.} Pour un univers $\U$ donné (contenant $\mathcal{N}$ comme élément), les conditions suivantes sont équivalentes :
\begin{enumerate}[label={}]
\item[1)] Il existe une catégorie $\U$-accessible $\cM$ ayant une sous-catégorie fermée non accessible % [sic] 1) -> a), la démonstration renvoie à a) et b)
\msp{page 484}
dans $\cM$.
\item[b)] Il existe une catégorie $\cC \in \U$, et une suite transfinie d'éléments de
\[
\cC^\wedge \ \overset{\text{\upshape déf}}{=}\ \underline{\mathrm{Hom}}(\cC, \mathrm{Ens}_{\U}) \ ,
\]
$\underline{x} = (x_\alpha)_{\alpha\in I}$, indexée par l'ensemble de tous les ordinaux $\alpha \in \U$, et telle que
\[
\Hom(x_\alpha, x_\beta) \ =\ \emptyset \qquad\qquad \textit{si } \alpha < \beta.
\]
\end{enumerate}
\endgroup

\textsc{Démonstration}. On a b) $\Longrightarrow$ a) en prenant $\cM = \cC^\wedge$, en vertu de la proposition précédente. On voit de même a) $\Longrightarrow$ b), en plongeant $\cM$ dans une catégorie $\cC^\wedge$ ($\cC$ étant une petite sous-catégorie strictement génératrice de $\cM$).

Je n'ai pas été capable (\footnote{Mais si, finalement, cf. 4.16.7.10 plus bas.}) de construire une suite $(x_\alpha)_{\alpha\in I}$ comme exigée dans cette proposition, ni de prouver qu'une telle suite n'existe pas, même dans un cas particulier tel que $\cM = (k\text{-alg. comm.})$, $k$ un anneau commutatif. Je n'ai pas
\msp{page 485}
même prouvé qu'on ne puisse trouver une suite transfinie de \textit{corps} $k_\alpha$, ayant les propriétés dites.

A fortiori, je n'ai pu tirer au clair la question si une sous-catégorie $\underline{\mathrm{L}}$-stable d'une catégorie $\cM$ accessible est toujours accessible dans $\cM$ (\footnote{C'est faux, cf. plus bas.}). Il se pourrait qu'il n'y ait pas de tel cas. Il se pourrait aussi qu'il y en ait, même avec une sous-catégorie fermée (laquelle est stable par tout type de $\varinjlim$ représentable dans $\cM$, sauf la $\varinjlim$ sur une catégorie d'indices vide $\ldots$).

\printedpage{188}
% head: A. Grothendieck : Dérivateurs, Ch. XVIII, Édition des Universités de Montpellier II et Paris VII
% foot: Version 17 décembre 2025

Pour les sous-catégories \textit{ouvertes}, je rappelle que j'ai construit (dans l'appendice) un exemple d'une telle sous-catégorie, non accessible dans $\cM$. Au sujet des sous-catégories ouvertes, je signale ici l'énoncé (quasi tautologique) :

\msp{page 486}
\textbf{Proposition 4.16.7.6} (\footnote{Cette proposition et son corollaire sont contenus dans 4.16.7.8 ci-dessous, il n'y a pas lieu d'en faire un énoncé séparé.}). \textit{Soient $\cM$ une catégorie accessible, $\uS$ une sous-catégorie \textup{ouverte} de $\cM$. Pour que $\uS$ soit accessible dans $\cM$, il faut et il suffit qu'elle soit $\underline{L}$-stable dans $\cM$, i.e. qu'il existe un cardinal $\pi_0$ tel que $\cM$ satisfasse à $L_{\pi_0}$, et que $\uS$ soit stable dans $\cM$ par $\varinjlim$ de type $L_{\pi_0}$. De façon plus précise, si cette condition est satisfaite, alors pour tout cardinal $\pi$ tel que $\pi \geqslant \pi_0$ et que $\pi$ soit utile pour $\cM$, $\uS$ est $\pi$-accessible dans $\cM$. Ou encore : si $\pi$ est un cardinal utile pour $\cM$, alors $\uS$ est $\pi$-accessible dans $\cM$ si et seulement s'il est $L_\pi$-stable dans $\cM$.}

Cette dernière assertion résulte aussitôt du corollaire suivant, lui-même tautologique.

\textbf{Corollaire 4.16.7.7.} \textit{Soit $\pi$ un cardinal comme indiqué ci-dessus, soit $\uS(\pi) = \uS \cap \cM_\pi$, et soit $x \in \Ob \cM$ et
\[
x \ \simeq\ \varinjlim_{i\in I} x_i \qquad\qquad
\begin{cases}
x_i \in \Ob \cM_\pi \\[1ex]
I \ \textit{grande devant}\ \pi
\end{cases}
\]
une représentation de $x$ comme image essentielle d'un objet de $\Ind_\pi \cM_\pi$. Alors $x \in \Ob \uS$ $\Longleftrightarrow$ $x_i \in \Ob \uS(\pi)$ $\forall\, i \in I$.}

\msp{page 487}
J'ai finalement tiré au clair la question de l'accessibilité des sous-catégories ouvertes et des sous-catégories fermées d'une catégorie accessible :

\begingroup\itshape
\noindent\textbf{\upshape Proposition 4.16.7.8.} Soient $\cM$ une catégorie accessible, $\uS$ une sous-catégorie fermée de $\cM$, $\uT$ la sous-catégorie ouverte complémentaire de $\uS$, $\pi$ un cardinal infini utile pour $\cM$. Les conditions suivantes sont équivalentes :
\begin{enumerate}[label={}]
\item[a)] $\uS$ est $\pi$-accessible dans $\cM$.
\item[a$'$)] $\uS$ est le fermé engendré par $\uS(\pi) \overset{\text{\upshape déf}}{=} \uS \cap \cM_\pi$, i.e. pour $x \in \Ob \cM$, on a $x \in \Ob \uS$ si et seulement si $\exists\, \xi \in \Ob \uS(\pi)$ tel que $\Hom_{\cM}(\xi, x) \neq \emptyset$.
\item[b)] $\uT$ est $\pi$-accessible dans $\cM$.
\item[b$'$)] $\uT$ est $\underline{L}_\pi$-stable dans $\cM$.
\end{enumerate}
\endgroup

\noindent Notons tout de suite le

\begingroup\itshape
\noindent\textbf{\upshape Corollaire 4.16.7.9.} Pour $\cM$, $\uS$, $\uT$ comme dessus, conditions équivalentes :
\begin{enumerate}[label={}]
\item[a)] $\uS$ est accessible dans $\cM$.
\item[a$'$)] Il existe une \textup{petite} partie $A$ de $\Ob \cM$, telle que $\uS$ soit la sous-catégorie fermée ${}_A\cM$ engendrée par $A$, i.e. $\uS = \bigcup_{\xi\in A} {}_\xi\cM$, où ${}_\xi\cM$ est la sous-catégorie fermée de $\cM$ formée des $x \in \Ob \cM$ tel que $\Hom(\xi, x) \neq \emptyset$.
\item[b)] \msp{page 488}$\uT$ est accessible dans $\cM$.
\end{enumerate}
\endgroup

\end{document}
